% Options for packages loaded elsewhere
\PassOptionsToPackage{unicode}{hyperref}
\PassOptionsToPackage{hyphens}{url}
%
\documentclass[
]{article}
\usepackage{lmodern}
\usepackage{amssymb,amsmath}
\usepackage{ifxetex,ifluatex}
\ifnum 0\ifxetex 1\fi\ifluatex 1\fi=0 % if pdftex
  \usepackage[T1]{fontenc}
  \usepackage[utf8]{inputenc}
  \usepackage{textcomp} % provide euro and other symbols
\else % if luatex or xetex
  \usepackage{unicode-math}
  \defaultfontfeatures{Scale=MatchLowercase}
  \defaultfontfeatures[\rmfamily]{Ligatures=TeX,Scale=1}
\fi
% Use upquote if available, for straight quotes in verbatim environments
\IfFileExists{upquote.sty}{\usepackage{upquote}}{}
\IfFileExists{microtype.sty}{% use microtype if available
  \usepackage[]{microtype}
  \UseMicrotypeSet[protrusion]{basicmath} % disable protrusion for tt fonts
}{}
\makeatletter
\@ifundefined{KOMAClassName}{% if non-KOMA class
  \IfFileExists{parskip.sty}{%
    \usepackage{parskip}
  }{% else
    \setlength{\parindent}{0pt}
    \setlength{\parskip}{6pt plus 2pt minus 1pt}}
}{% if KOMA class
  \KOMAoptions{parskip=half}}
\makeatother
\usepackage{xcolor}
\IfFileExists{xurl.sty}{\usepackage{xurl}}{} % add URL line breaks if available
\IfFileExists{bookmark.sty}{\usepackage{bookmark}}{\usepackage{hyperref}}
\hypersetup{
  pdftitle={Learning - pilot},
  pdfauthor={Matan Mazor},
  hidelinks,
  pdfcreator={LaTeX via pandoc}}
\urlstyle{same} % disable monospaced font for URLs
\usepackage[margin=1in]{geometry}
\usepackage{graphicx,grffile}
\makeatletter
\def\maxwidth{\ifdim\Gin@nat@width>\linewidth\linewidth\else\Gin@nat@width\fi}
\def\maxheight{\ifdim\Gin@nat@height>\textheight\textheight\else\Gin@nat@height\fi}
\makeatother
% Scale images if necessary, so that they will not overflow the page
% margins by default, and it is still possible to overwrite the defaults
% using explicit options in \includegraphics[width, height, ...]{}
\setkeys{Gin}{width=\maxwidth,height=\maxheight,keepaspectratio}
% Set default figure placement to htbp
\makeatletter
\def\fps@figure{htbp}
\makeatother
\setlength{\emergencystretch}{3em} % prevent overfull lines
\providecommand{\tightlist}{%
  \setlength{\itemsep}{0pt}\setlength{\parskip}{0pt}}
\setcounter{secnumdepth}{-\maxdimen} % remove section numbering

\title{Learning - pilot}
\author{Matan Mazor}
\date{06/07/2020}

\begin{document}
\maketitle

\hypertarget{methods}{%
\section{Methods}\label{methods}}

\hypertarget{participants}{%
\subsection{Participants}\label{participants}}

A total of 22 participants were recruited from Prolific. Participants
were Native English Speakers and had a minimum approval rate of 0.9.
They were paid £0.80 for their participation, in addition to a bonus
payment of up to £0.64. The experiment took 7 minutes to complete.

\hypertarget{experimental-design}{%
\subsection{Experimental design}\label{experimental-design}}

On each trial, participants were presented with one or two black shapes,
presented inside a frame in the center of the screen. They then used
their keyboard to indicate whether they thought a cat will appear after
the shape (`J' key) or not (`F' key). Following their response, the
shapes disappeared from the screen. On half of the trials, the frame
remained empty for 2 seconds. On the other half, an image of a cat
appeared inside the frame. Participants were awarded a point for
accurate predictions, and points were later translated to a bonus
payment with a conversion rate of £0.01 per point. Their total number of
point and trial count remained visible on the screen throughout the
experiment.

After Reading the information sheet and signing the consent form,
participants were instructed about the task. In the following practice
part, participants learned to associate one shape with the appearance of
a cat, and one shape with the absence of a cat. Shapes were randomly
selected separately for each participant from the following: circle,
hexagon, up-facing triangle, down-facing triangle, 5-armed star, heart,
spade, club, diamond, and cross. Once they performed with 100\% accuracy
for a block of four trials, they moved on to a verbal comprehension test
(`Please describe the task in your own words'), and then moved on to the
main task.

The main task consisted of 64 trials, separated into two parts of 32
trials each. Trials were presented in 8 blocks, each block comprising 8
different trial types, as follows:

\begin{enumerate}
\def\labelenumi{\arabic{enumi}.}
\tightlist
\item
  Stimulus A, coupled with cat presence.
\item
  Stimulus B, coupled with cat absence.
\item
  Stimulus C, coupled with cat presence,
\item
  Stimuli C and D, coupled with cat absence.
\item
  Stimuli E and F, coupled with cat presence.
\item
  Stimulus E, coupled with cat absence.
\end{enumerate}

Stimuli A-F were randomly sampled, without replacement, from the
remaining shapes that were not used as example stimuli in the practice
phase.

For 14 of the participants (Q positive group), the two remaining
trial-types were:

\begin{enumerate}
\def\labelenumi{\arabic{enumi}.}
\setcounter{enumi}{6}
\tightlist
\item
  A ring intersected by a line (Q), coupled with cat presence.
\item
  A ring (O), coupled with cat absence.
\end{enumerate}

For the remaining 8 participants (Q negative group), the two remaining
trial-types were:

\begin{enumerate}
\def\labelenumi{\arabic{enumi}.}
\setcounter{enumi}{6}
\tightlist
\item
  A ring (O), coupled with cat presence.
\item
  A ring intersected by a ling (Q), coupled with cat absence.
\end{enumerate}

Trial types 1 and 2 formed the symmetric learning condition, where one
stimulus predicted the presence of a cat, and one stimulus predicted its
absence. Trial types 3 and 4 formed the Feature-Negative condition,
where the presence of a feature (stimulus D) predicted the absence of a
cat. Trial types 4 and 5 formed the Feature-Positive condition, where
the presence of a feature (stimulus F) predicted the presence of a cat.
Finally, trial types 7 and 8 formed the test case, where we tested
whether the assignment of stimuli to the presence or absence of a cat
makes learning more or less efficient.

\hypertarget{exclusion}{%
\subsection{Exclusion}\label{exclusion}}

Participants were excluded for responding too quickly (RT\textless250)
in more than 14 trials. After exclusion, we had datasets from 6
participants in the Q positive and 8 in the Q negative groups,
respectively.

\hypertarget{results}{%
\subsection{Results}\label{results}}

Overall accuracy was high (\(M = 0.74\), 95\% CI \([0.67\), \(0.81]\)).
Analysis of variance was applied to the data from trial-types 1-6 (not
including the test stimuli), with the dependent variable accuracy and
independent variables block number, stimulus pair, and participant. As
expected, block number had a significant effect on accuracy, suggesting
that participants learned to associate stimuli with outcome over the
course of the experiment (\(F(1, 319) = 85.20\),
\(\mathit{MSE} = 0.08\), \(p < .001\)). Furthermore, accuracy was
significantly different for the different stimulus pairs
(\(F(2, 319) = 9.19\), \(\mathit{MSE} = 0.08\), \(p < .001\); see Fig.
\ref{fig:learning_curves}).

However, this difference was driven by faster learning for the symmetric
learning condition (trial types 1 and 2) compared to the other learning
conditions (\(M_d = 0.13\), 95\% CI \([0.04\), \(0.23]\),
\(t(13) = 3.13\), \(p = .008\)). No significance difference was observed
in overall accuracy between the feature-positive and feature-negative
learning accuracies (\(M_d = 0.06\), 95\% CI \([-0.04\), \(0.15]\),
\(t(13) = 1.26\), \(p = .228\)).

\begin{figure}
\centering
\includegraphics{Lotz_files/figure-latex/fig:learning_curves-1.pdf}
\caption{Learning curves for the three within-participant conditions}
\end{figure}

\begin{figure}
\centering
\includegraphics{Lotz_files/figure-latex/fig:RT_curves-1.pdf}
\caption{Learning curves for the three within-participant conditions}
\end{figure}

\end{document}
